
\subsection*{1.4}
\phantomsection
\label{sga3:VIIB:1.4}
\addcontentsline{toc}{subsection}{1.4. Descent and exactness}

\noindent\textbf{Theorem 1.4.}
\phantomsection
\label{sga3:VIIB:1.4.thm}
{\itshape Let \(k\) be a pseudocompact ring and let
\[
  d_0,d_1:X_1\rightrightarrows X
\]
be an equivalence pair in \(\operatorname{Vaf}/k\)
\textup{(cf.\ Expos\'e V, \S~2.b)} such that \(d_1\) is
topologically flat.}
\begin{enumerate}
  \item[\textup{(i)}]
  {\itshape The canonical projection from \(X\) onto
  \(X/X_1=\operatorname{Coker}(d_0,d_1)\) is surjective and
  topologically flat, and the morphism
  \[
    X_1\longrightarrow X\times_{X/X_1}X
  \]
  with components \(d_0\) and \(d_1\) is an isomorphism.}

  \item[\textup{(ii)}]
  {\itshape If \(X\) is topologically flat over \(k\), then so is
  \(X/X_1\).}
\end{enumerate}

First note that \textup{\SGARefLink{sga3:vii:paragraph:5-1:equation:ii}{(ii)}} follows from \textup{\SGARefLink{sga3:vii:paragraph:5-1:equation:i}{(i)}} by
\SGARefLink{sga3:VIIB:1.3.3}{1.3.3}\,\textup{\SGARefLink{sga3:vii:paragraph:5-1:equation:ii}{(ii)}}. The proof of \textup{\SGARefLink{sga3:vii:paragraph:5-1:equation:i}{(i)}} occupies
\SGARefLink{sga3:VIIA:1.4.1}{1.4.1}, \SGARefLink{sga3:VIIA:1.4.2}{1.4.2}, and \SGARefLink{sga3:VIIB:1.4.3}{1.4.3}.

\subsubsection*{1.4.1}
\phantomsection
\label{sga3:VIIB:1.4.1}

We first show that we may reduce to the case in which \(X\) has only
one point. Since we are dealing with an equivalence pair, one sees,
as in Expos\'e V, \S~3.e), that an equivalence relation is defined on
the underlying set of \(X\) by declaring two points \(x,y\) to be
equivalent if there exists a point \(z\in X_1\) such that
\[
  d_0(z)=y
  \qquad\text{and}\qquad
  d_1(z)=x.
\]
We may plainly suppose, without loss of generality, that \(X\)
contains a single equivalence class for this relation; in other words,
that \(X/X_1\) has only one point \textup{(see the construction of
\(X/X_1\) given in \SGARefLink{sga3:VIIB:1.2}{1.2})}.

In this case, let \(x\) be any point of \(X\), and let \(U\) be the
formal variety having \(x\) as its only point and the same local ring
as \(X\) at \(x\). As in Expos\'e V, \S~6, one then sees that the
equivalence relation induced by \((d_0,d_1)\) on \(U\) still satisfies
the hypotheses of the theorem and that it is enough to prove the
theorem for this latter equivalence relation \textup{(\(U\) is a
``quasi-section'')}.

Let us briefly recall the principle of the proof given in
Expos\'e V, \S~6. Put
\[
  V=d_0^{-1}(U)=U\times_{i,d_0}X_1,
\]
where \(i\) is the inclusion of \(U\) into \(X\), and let \(u\) and
\(v\) be the morphisms with source \(V\) induced respectively by
\(d_0\) and \(d_1\):
\begin{center}
  \phantomsection
  \label{sga3:VIIB:1.4.1:quasisection-span}
  \(
X\xleftarrow{\ v\ }V\xrightarrow{\ u\ }U.
\)

\end{center}
It is clear that \(u\) and \(v\) are topologically flat and that
\(u\) is surjective; since \(X\) contains a single equivalence class,
\(v\) is surjective. If \((v_0,v_1)\) is the inverse image of the
equivalence pair \((d_0,d_1)\) under \(v\) \textup{(cf.\ V, 3.a))},
it follows from V, 3.c) and 3.d) that \(X/X_1\) and the quotient of
\(U\) by the equivalence relation induced by \((d_0,d_1)\) both
identify with \(\operatorname{Coker}(v_0,v_1)\). As in the proof of
V, \SGARefLink{lem:sga3-V-6.1}{6.1}, one then sees that if the conclusion of \SGARefLink{sga3:VIIB:1.4.thm}{Theorem~1.4} holds for
\(U\), it also holds for \(X\).

\subsubsection*{1.4.2}
\phantomsection
\label{sga3:VIIB:1.4.2}

We are thus reduced to the case in which \(X\) has only one
point.
Consider the following commutative diagram
\textup{(cf.\ V, \S~1, \textup{(0,1,2)})}:
\begin{center}
  \phantomsection
  \label{sga3:VIIB:1.4.2:equivalence-groupoid-diagram}
  \begin{tikzcd}[row sep=3.3em,column sep=3.9em]
X_2
  \arrow[r,shift left=.6ex,"d'_1"]
  \arrow[r,shift right=.6ex,"d'_0"']
  \arrow[d,"d'_2"']
  & X_1 \arrow[r,"d_0"] \arrow[d,"d_1"]
  & X \arrow[d] \\
X_1
  \arrow[r,shift left=.6ex,"d_1"]
  \arrow[r,shift right=.6ex,"d_0"']
  & X \arrow[r]
  & X/X_1
\end{tikzcd}

\end{center}
Here
\[
  X_2=X_1\times_{d_1,d_0}X_1,
\]
and \(d'_0,d'_1,d'_2\) are respectively the
morphisms
\[
  (x,y,z)\longmapsto(x,y),\qquad
  (x,y,z)\longmapsto(x,z),\qquad
  (x,y,z)\longmapsto(y,z).
\]

If \(B,A,A_1,A_2\) denote respectively the affine algebras of
\(X/X_1,X,X_1,X_2\), the preceding diagram induces the commutative
diagram
\begin{center}
  \phantomsection
  \label{sga3:VIIB:1.4.2:affine-algebra-diagram}
  \begin{tikzcd}[row sep=3.3em,column sep=4.1em]
A_2
  & A_1
      \arrow[l,shift left=.6ex,"j_0"]
      \arrow[l,shift right=.6ex,"j_1"']
  & A \arrow[l,"i_0"'] \\
A_1 \arrow[u,"j_2"]
  & A
      \arrow[l,shift left=.6ex,"i_0"]
      \arrow[l,shift right=.6ex,"i_1"']
      \arrow[u,"i_1"']
  & B \arrow[l,"i"'] \arrow[u,"i"]
\end{tikzcd}

\end{center}
in which both rows are exact and the squares determined by
\((i_0,j_0)\) and \((i_1,j_1)\) are cocartesian. Since the morphism
\[
  X_1\longrightarrow X\times X
\]
with components \(d_0\) and \(d_1\) is a monomorphism by hypothesis,
the morphism
\[
  A\widehat{\otimes}_k A\longrightarrow A_1
\]
with components \(i_0\) and \(i_1\) is surjective by
\SGARefLink{sga3:VIIB:1.3.prop}{Proposition~1.3}.

\begin{samepage}
This means that \(i_1\) makes \(A_1\) a pseudocompact \(A\)-module,
topologically free by hypothesis, generated by \(i_0(A)\). Since
\(A\) is local, \SGARefLink{sga3:VIIB:1.4.3}{Lemma~1.4.3} below gives a topologically free
\(k\)-module \(V\) and a morphism of pseudocompact \(k\)-modules
\(f:V\to A\) such that the morphism
\[
  \alpha_1:
  A\widehat{\otimes}_kV\longrightarrow A_1,
  \qquad
  a\widehat{\otimes}v
  \longmapsto
  i_1(a)\cdot i_0\bigl(f(v)\bigr)
\]
is an isomorphism.
\end{samepage}
\begin{samepage}
Let
\[
  \alpha:B\widehat{\otimes}_kV\longrightarrow A
  \qquad\text{and}\qquad
  \alpha_2:A_1\widehat{\otimes}_kV\longrightarrow A_2
\]
be the morphisms
\[
  b\widehat{\otimes}v
  \longmapsto
  i(b)\cdot f(v)
  \qquad\text{and}\qquad
  a_1\widehat{\otimes}v
  \longmapsto
  j_2(a_1)\cdot j_0i_0\bigl(f(v)\bigr).
\]
\end{samepage}
In the following commutative diagram, both rows are exact and the two
left-hand squares are cocartesian. Since \(\alpha_1\) is an
isomorphism, the same is true of \(\alpha_2\), and hence of
\(\alpha\).
\begin{center}
  \phantomsection
  \label{sga3:VIIB:1.4.2:alpha-cocartesian-diagram}
  \begin{tikzcd}[row sep=3.4em,column sep=3.5em]
A_2
  & A_1
      \arrow[l,shift left=.6ex,"j_0"]
      \arrow[l,shift right=.6ex,"j_1"']
  & A \arrow[l,"i_0"'] \\
A_1\widehat\otimes_k V \arrow[u,"\alpha_2"]
  & A\widehat\otimes_k V
      \arrow[l,shift left=.6ex,"{i_0\widehat\otimes V}"]
      \arrow[l,shift right=.6ex,"{i_1\widehat\otimes V}"']
      \arrow[u,"\alpha_1"']
  & B\widehat\otimes_k V
      \arrow[l,"{i\widehat\otimes V}"']
      \arrow[u,"\alpha"']
      \, .
\end{tikzcd}

\end{center}

This shows, on the one hand, that \(A\) is topologically free over
\(B\), with pseudobasis \(f(V)\) \textup{(cf.\ \SGARefLink{sga3:VIIB:0.2.1}{0.2.1})}, and that a
pseudobasis of \(A_1\) over \(A\), where \(A_1\) is regarded as an
\(A\)-module through \(i_1\), is obtained by taking the image under
\(i_0\) of \(f(V)\). It follows that the morphism
\[
  A\widehat{\otimes}_B A\longrightarrow A_1
\]
with components \(i_1\) and \(i_0\) is an isomorphism:
\[
  A\widehat{\otimes}_B A
  \simeq
  A\widehat{\otimes}_B B\widehat{\otimes}_kV
  \simeq
  A\widehat{\otimes}_kV
  \simeq
  A_1.
\]
This proves \SGARefLink{sga3:VIIB:1.4.thm}{Theorem~1.4}, modulo the following lemma.

\subsubsection*{1.4.3}
\phantomsection
\label{sga3:VIIB:1.4.3}

\begin{samepage}
\noindent\textbf{Lemma 1.4.3.}
{\itshape Let \(k\) be a pseudocompact ring, let \(A\) be a local
profinite \(k\)-algebra, let \(A_1\) be a topologically free
\(A\)-module, and let
\[
  i_0:M\longrightarrow A_1
\]
be a morphism of pseudocompact \(k\)-modules. Suppose that the map
\[
  A\widehat{\otimes}_kM\longrightarrow A_1,
  \qquad
  a\widehat{\otimes}m
  \longmapsto
  a\cdot i_0(m),
\]
is surjective. Then there exist a topologically free \(k\)-module
\(V\) and a morphism of pseudocompact \(k\)-modules
\[
  f:V\longrightarrow M
\]
such that the map
\[
  A\widehat{\otimes}_kV\longrightarrow A_1,
  \qquad
  a\widehat{\otimes}v
  \longmapsto
  a\cdot i_0\bigl(f(v)\bigr),
\]
is bijective.}
\end{samepage}

Since every pseudocompact \(k\)-module is a quotient of a
topologically free \(k\)-module, we may suppose, without loss of generality, that
\(M\) is topologically free. Thus write
\[
  M=\prod_{i\in I}M_i,
\]
where each \(M_i\) is a copy of \(k\). In this case,
\[
  A\widehat{\otimes}_kM
  =
  \prod_{i\in I}
  A\widehat{\otimes}_kM_i.
\]
Since the map
\[
  a\widehat{\otimes}m\longmapsto a\cdot i_0(m)
\]
is surjective and \(A_1\) is projective, its kernel is a direct
factor of \(A\widehat{\otimes}_kM\). As \(A\) is local, the exchange
theorem \textup{(\SGARefLink{sga3:VIIB:0.3.4}{0.3.4})} shows that this kernel has as a complement a
partial product
\[
  \prod_{i\in J}A\widehat{\otimes}_kM_i,
\]
where \(J\) is a subset of \(I\). One may therefore take
\[
  V=\prod_{i\in J}M_i.
\]

